\section{The polylogarithmic boundary estimate and the unit extension}
\label{sec:faithful-boundary-unit}
We follow the multiplication-and-division proof of logarithmic
removability in \srcid{ContinuousNormMulOneAddAbsLogPow}, followed
by the actual subframe and Hodge Gram comparisons in
\srcid{BoundaryChartFrameGrowth}. The construction of the period
family used in those comparisons remains a separate input.

For a finite coordinate set \(S\subset\{1,\ldots,n\}\), put
\[
 D_S=\{z:\prod_{i\in S}z_i=0\},\quad
 P_S(z)=\prod_{i\in S}z_i,\quad
 \ell_S(z)=1+\sum_{i\in S}|\log|z_i||\quad(z\notin D_S).
 \tag{BG1}\label{eq:faithful-log-weight}
\]
Local logarithmic growth means that, on each compact subset of
an open \(U\subset\mathbb C^n\), the norm off \(D_S\) is bounded
by \(C\ell_S^N\), for some \(C\ge1\) and integer \(N\ge0\).

\begin{srlemma}[Removal of logarithmic growth]
\label{lem:faithful-removable-log-growth}
A holomorphic function on \(U\setminus D_S\) with local logarithmic
growth has a unique holomorphic extension to \(U\).
\end{srlemma}
\begin{proof}
Give \(f\) any finite values on \(D_S\) when writing products below.
For every integer \(a\ge1\),
\(r(1+|\log r|)^a\to0\) as \(r\downarrow0\):
for \(r\le1\), set \(x=1-\log r\), obtaining \(e\,x^ae^{-x}\).
Extending its value at zero by zero makes this a continuous
function on \([0,\infty)\).
For nonnegative \(a_i\), induction gives
\(1+\sum a_i\le\prod(1+a_i)\), since the mixed terms are nonnegative.
On a compact neighborhood with \(|f|\le C\ell_S^N\) this yields
\[
 |P_Sf|\le C\prod_{i\in S}
       \{|z_i|(1+|\log|z_i||)^{N+1}\}\quad(z\notin D_S).
 \tag{BG2}\label{eq:faithful-product-log-majorant}
\]
Each factor has a continuous finite limit, and at a divisor
point at least one has limit zero.
Since \(P_Sf=0\) on the divisor, it follows that \(P_Sf\to0\)
at every divisor point.
This is \srcid{HasLogGrowth.tendsto_prod_mul}.

Set \(G=P_S^2f\) off the divisor and zero on it.
At \(p\in D_S\), the polynomial \(P_S\) is
\(O(\|z-p\|)\), while \(P_Sf=o(1)\).
Thus \(G=o(\|z-p\|)\).
It is complex Fr\'echet differentiable at \(p\), with derivative zero.
Away from the divisor it is holomorphic, so \(G\) is holomorphic on \(U\).

If a holomorphic \(H\) vanishes on \(\{z_i=0\}\cap U\), it has a
holomorphic quotient \(H_1\) with \(H=z_iH_1\).
Locally, the convergent power series of \(H\) in the \(i\)-th
coordinate has zero constant coefficient; shifting its powers
down by one gives the quotient.
Off the hyperplane it is \(H/z_i\), and on the hyperplane its
value is \(\partial_iH\).
These descriptions agree and glue.
This is \srcid{exists_eq_coord_mul_of_isOpen}.

Suppose more generally that \(H=z_if_1\) off \(D_S\), with
\(i\in S\) and \(f_1\) of local logarithmic growth.
To prove \(H(p)=0\) for every \(p_i=0\), replace all zero
coordinates of \(p\) by the same nonzero \(t\), leaving the
others fixed; call this point \(\gamma(t)\).
For small \(t\), it lies outside \(D_S\) in a fixed compact
neighborhood in \(U\), and for some fixed \(M\ge1\)
\[
 \ell_S(\gamma(t))\le M(1+|\log|t||),\qquad
 |H(\gamma(t))|\le |t|CM^N(1+|\log|t||)^N\longrightarrow0.
\]
Continuity proves the required vanishing, including intersections
of hyperplanes.

List the coordinates of \(S\) twice.
Starting from \(G=P_S^2f\), divide successively by the coordinates
in this list.
At each step the remaining coordinate product times \(f\)
has local logarithmic growth, since a polynomial is bounded
on compact subsets.
The preceding paragraph proves vanishing on the next hyperplane,
so holomorphic division applies.
Cancellation off \(D_S\) preserves equality with the remaining
product times \(f\).
After the doubled list is exhausted, the quotient is the
required extension.
This is \srcid{HasLogGrowth.exists_extension_of_prod_mul}.
Uniqueness follows from continuity and density of the complement,
also witnessed by the paths \(\gamma(t)\).
For empty \(S\) no division is necessary.
\end{proof}

\begin{srcorollary}[Two-sided growth gives a unit]
\label{cor:faithful-log-growth-unit}
If \(f\) is holomorphic and nonzero off \(D_S\), and both \(f\)
and \(f^{-1}\) have local logarithmic growth, its holomorphic
extension is nowhere zero.
\end{srcorollary}
\begin{proof}
Extend \(f\) and \(f^{-1}\) separately by the lemma.
Their product is \(1\) on the dense complement and hence everywhere.
This is \srcid{exists_unit_extension_of_log_growth}.
A bound on \(f\) alone would not establish nonvanishing.
\end{proof}

\subsection{The subframe determinant estimate}
Matrix inequalities below mean nonnegativity of Hermitian quadratic
forms. For a complex \(g\times k\) matrix \(P\) and positive real
symmetric \(g\times g\) matrix \(Y\), write
\[
 G(P)=\det(P^t\overline P),\qquad
 D(P,Y)=\det(P^tY^{-1}\overline P).
 \tag{BG3}\label{eq:faithful-subframe-determinants}
\]
For full-rank \(P\) these are positive real numbers.
The conjugations agree with the source: use the columns of
\(\overline P\) in the usual Hermitian Gram formula.

\begin{srlemma}[Inverse and determinant comparison]
\label{lem:faithful-subframe-determinant-comparison}
Let \(P\) have full column rank. If \(0<cI\le Y\le CrI\)
and \(r\ge1\), then
\[
              (Cr)^{-k}G(P)\le D(P,Y)\le c^{-k}G(P).
 \tag{BG4}\label{eq:faithful-subframe-det-bounds}
\]
\end{srlemma}
\begin{proof}
For \(g>0\) the matrix inequalities imply \(C>0\).
Orthogonal diagonalization gives
\((Cr)^{-1}I\le Y^{-1}\le c^{-1}I\).
They hold also on complex vectors: for \(u=a+ib\),
\(u^*Y^{-1}u=a^tY^{-1}a+b^tY^{-1}b\), since symmetry
cancels the imaginary cross terms.
Conjugating by the injective column map \(\overline P\) gives
\[
 (Cr)^{-1}(P^t\overline P)
 \le P^tY^{-1}\overline P\le c^{-1}(P^t\overline P).
\]
Conjugate by the inverse positive square root of
\(P^t\overline P\). The middle matrix has its \(k\) positive
eigenvalues between \((Cr)^{-1}\) and \(c^{-1}\).
Multiplying them proves the determinant inequalities.
These are the source's
\srcid{subframe_inv_bounds} and \srcid{subframe_det_bounds}.
If \(g=0\), full rank forces \(k=0\) and both determinants are \(1\).
\end{proof}

\begin{srproposition}[The compact-family bound]
\label{prop:faithful-compact-subframe-log}
Let \(P(x)\) be continuous of rank \(k\) on a compact set \(T\).
Fix \(c>0\) and real \(C\).
There is \(C'>0\), independent of \(x,Y,r\), such that
\(0<cI\le Y\le CrI\) and \(r\ge1\) imply
\[
                 D(P(x),Y)\le C',\qquad
                    1\le C'r^kD(P(x),Y).
 \tag{BG5}\label{eq:faithful-compact-subframe-log}
\]
\end{srproposition}
\begin{proof}
If \(T\) is empty take \(C'=1\).
Otherwise \(G(P(x))\) is positive and continuous; let its
minimum be \(\mu>0\), and let \(M\) bound its absolute value.
The source takes
\[
                        C'=c^{-k}|M|+|C|^k/\mu+1.
 \tag{BG6}\label{eq:faithful-subframe-constant}
\]
For \(g>0\), (BG4) bounds the determinant above by the first
summand. Its lower bound gives
\(1\le(|C|^k/\mu)r^kD(P(x),Y)\), proving the other assertion.
For \(g=k=0\), both determinants are \(1\) and \(C'\ge1\).
This proves \srcid{subframe_det_log_growth}.
Uniform full rank on the compact set is essential for \(\mu>0\).
\end{proof}

For the source's logarithmic period datum, a holomorphic matrix
\(A(z)\), invertible at \(0\), and a period matrix \(\tau(w)\)
on the logarithmic cover satisfy
\[
 A(z)H(z)A(z)^*=(\operatorname{Im}\tau(w))^{-1},
                         \qquad z=\exp(w).
 \tag{BG7}\label{eq:faithful-period-trivialization}
\]
The period order bounds are
\[
 cI\le Y=\operatorname{Im}\tau(w)\le C R(z)I,\qquad
 R(z)=e+2\sum_{i<a}|\log|z_i||,\qquad
       1\le R(z)\le(e+2)\ell_S(z).
 \tag{BG8}\label{eq:faithful-chart-radius}
\]
The source supplies these bounds through \srcid{ImTauBounds.cond}
with the \srcid{RiemannExtensionSlices},
\srcid{PuncturedExpOrder}, and \srcid{LogCoverDescent} inputs.
We retain this precise period-family input.

Let \(F(z)\) be continuous, of row rank \(k\) at \(0\).
After shrinking, the matrices \((F(z)A(z)^{-1})^t\) lie in one
compact neighborhood of full-rank matrices.
Indeed their Euclidean Gram determinant is positive at \(0\)
and continuous; choose a compact neighborhood inside its positive
locus, and use continuity of \(F\) and matrix inversion.
Apply the preceding proposition and (BG7):
\[
 FHF^*=(FA^{-1})Y^{-1}(FA^{-1})^*,\qquad
 |\det(FHF^*)|\le K,\quad
       1\le K\ell_S^k|\det(FHF^*)|,
 \tag{BG9}\label{eq:faithful-log-period-subframe}
\]
where \(K=C'(e+2)^k\).
The determinant is nonnegative real, so its norm is its real part.
This is \srcid{logPeriods_subframe_det_bounds} with exactly
the period-trivialization and order inputs just stated.

\subsection{The actual Hodge Gram factors and the flat norm}
In a boundary chart write \(F(z)_{ib}=\psi_z(u_i(b))\)
for the chosen model tangent columns, independent at \(0\).
The logarithmic period construction, after a radial rescaling,
identifies their actual Hodge Gram matrix as
\[
                         \Gamma=\pi FHF^* .
 \tag{BG10}\label{eq:faithful-hodge-gram}
\]
For any \(a\ne0\), the triangle inequality for logarithms gives
\[
 \ell_S(az)\le\{1+n|\log|a||\}\ell_S(z).
 \tag{BG11}\label{eq:faithful-log-weight-rescaling}
\]
There are at most \(n\) terms in \(S\), and \(\ell_S\ge1\).
Thus, if the chart rescaling is \(z\mapsto\rho z\) with \(\rho>0\),
(BG9) gives
\[
 |\det\Gamma|\le K,\qquad
             1\le K\ell_S^k|\det\Gamma|,\qquad
 K=K_0\pi^k\{1+n|\log|\rho^{-1}||\}^k .
 \tag{BG12}\label{eq:faithful-hodge-gram-log-bound}
\]
Here \(K_0\) is the constant before this rescaling.
The factors \(\pi^k\) and the bracket are at least one;
multiplying by them preserves both estimates.
This is \srcid{hodgeGram_det_log_bounds}, including the factor
\(\pi\) and the radial change.

The two actual eigenframes of the flat analytic function have
Hodge Gram matrices \(\Gamma,\Delta\).
Its defining norm identity is
\[
                     |\Phi|^2|\det\Gamma|=|\det\Delta|.
 \tag{BG13}\label{eq:faithful-unitAn-gram-norm}
\]
Write \(x=|\Phi|>0\), \(g=|\det\Gamma|\), \(d=|\det\Delta|\).
The bounds for the same frames are
\[
 g\le K_1,\quad1\le K_1\ell_S^{N_1}g,\qquad
 d\le K_2,\quad1\le K_2\ell_S^{N_2}d.
\]
Multiplying the first lower bound by \(x^2\), then using
\(x^2g=d\le K_2\), proves
\(x^2\le K_1K_2\ell_S^{N_1}\).
Multiplying the second lower bound by \(x^{-2}\) similarly
proves \(x^{-2}\le K_1K_2\ell_S^{N_2}\).
Since \(y\le1+y^2\), \(\ell_S\ge1\), and the constants are
nonnegative, these imply the source's common bound
\[
           x,\ x^{-1}\le(1+|K_1K_2|)\ell_S^{N_1+N_2}.
 \tag{BG14}\label{eq:faithful-flat-two-sided-log}
\]
This is \srcid{norm_bounds_of_gram}.
No varying square-root or logarithm branch is chosen.

\begin{srproposition}[The boundary-chart unit]
\label{prop:faithful-bdry-chart-unit-extension}
Assume the actual boundary chart, its centre-independent tangent
columns, the analytic frame function's holomorphy and nonvanishing
clauses, its Gram identity (BG13), and the period inputs
(BG7)--(BG8) for those same tangent columns.
Its frame value extends to a holomorphic unit on a smaller whole chart.
\end{srproposition}
\begin{proof}
At a punctured point the actual abelian open contains a basic
neighborhood \(R[r^{-1}]\) with \(\psi_z(r)\ne0\).
Each value \(\psi_z(a)/\psi_z(r)^j\) is holomorphic where the
denominator is nonzero; this extends the coefficient maps
holomorphically to the localization.
The cartesian chart maps lift the model tangent columns over
this localization by their dual-number maps.
Their specializations are the stated eigenframes: their point maps
and both projections to the semiabelian pullbacks agree, so
cartesian uniqueness identifies the tangent vectors.
The holomorphy clause of the actual analytic frame function
therefore proves holomorphy of its displayed value.
This is the source's \srcid{differentiableAt_of_chart} argument.

Hodge Gram pullback identifies the two Gram matrices with these
same vectors. On the smaller of their two admissible radii,
(BG12)--(BG14) give logarithmic growth of the frame value and
its inverse. Nonvanishing on the punctured chart is part of
\srcid{IsUnitAn}.
Corollary~\ref{cor:faithful-log-growth-unit} gives the required unit.
This is \srcid{exists_unit_extension_of_bdryChart}.
\end{proof}

\paragraph{The remaining construction-specific period join.}
The actual logarithmic period datum is constructed in
\srcid{LogPeriodsRiemannForm.exists_logPeriods_isRiem}.
Its source proof builds the local family exponential, continues
the central fibre's discrete free periods, makes additive local
sections, identifies the Hodge matrix, and checks the torsion
pairing before producing the datum.
The change of period basis in
\srcid{PeriodMatrixLocalTrivialization} uses a symplectic marking
and \srcid{SiegelOfRiemann}.
The order bounds in (BG8) retain their precise source inputs.
Thus this section proves the determinant and removable-unit
arguments from those inputs; it does not infer their existence
from a monodromy slogan.
The actual chart and tangent-column comparison with the
resolution still has to be supplied by
\srcid{AffineChartBundleSections} before using
Section~\ref{sec:faithful-flat-line-extension}.

The frozen paths are
\srcid{src/Development/AbelianVarietyTheory/Proofs/CovolumeBounds/ContinuousNormMulOneAddAbsLogPow.lean};
\srcid{src/Development/AutomorphicSpectrum/Proofs/MetricProofs/ExistsUnitExtensionOfLogGrowth.lean},
\srcid{SubframeDetLogGrowth.lean} and
\srcid{LogPeriodsSubframeDetBounds.lean} in the same metric directory;
and
\srcid{src/Development/AutomorphicSpectrum/Proofs/GeometryProofs/HodgeGramBoundNearBoundary.lean}
and \srcid{BoundaryChartFrameGrowth.lean}.
All are read at \srcid{alpoge-source-7cbf7a0c}.
